\section{结论}
\label{sec:conclusion}

本文提出 \textbf{BudgetMem}，一个运行时智能体记忆框架，用于在按需记忆抽取中实现显式的性能--成本控制。
BudgetMem 为模块化记忆流水线中的每个模块配置 \textsc{低}/\textsc{中}/\textsc{高}预算档位，并学习一个轻量级路由器，在成本感知目标下选择档位。
以 BudgetMem 为统一试验平台，我们比较了三种互补的分层策略，即\textit{实现分层}、\textit{推理分层}与\textit{容量分层}，并刻画它们在不同预算下的权衡行为。
LoCoMo、LongMemEval 和 HotpotQA 上的实验表明，该方法在性能优先设置下具有强劲表现，并能在预算收紧时改善性能--成本前沿；实验还揭示了各策略分别在何种条件下带来最佳收益。
